\section{Introduction}
\label{sec:intro}

A large class of population-based optimizers is justified by a score. Candidate
solutions are placed on a graph, a per-neighbour quantity is computed---a
centrality of the graph, the quality of a neighbour's best---and the update is
weighted so that agents learn preferentially from the neighbours the score
distinguishes. The accompanying argument is causal: the method performs well
\emph{because} the score identifies informative agents. Where such mechanisms
are evaluated at all, it is by deletion.

Deletion cannot establish the claim. Removing a weighting changes the
optimizer's mixing operator, its effective neighbourhood and its convergence
rate, so a difference in outcome is consistent with the weights having been
merely unequal. The claim actually advanced is narrower: that the
\emph{assignment} of scores to nodes, and thereby to the solutions those nodes
hold, is informative. Testing it requires holding the mechanism at full
strength while destroying only that correspondence.

This paper supplies such a test, characterises it, and applies it at scale. The
first control permutes the scores across nodes: the score multiset, the weight
heterogeneity and the arithmetic cost are preserved exactly, while the pairing
of score to position is destroyed. Proposition~\ref{prop:mean} shows this
control has the uniform operator as its expectation while retaining the
original's weight dispersion in every realisation, so it interpolates precisely
between the two alternatives an ablation confounds. The second control replaces
an adaptive component's selection rule by a uniform draw over the same
alternatives on the same schedule, separating \emph{selecting} a component from
merely \emph{varying} it.

A control is informative only if the mechanism it neutralises was doing
something, and we make that precondition quantitative.
Proposition~\ref{prop:leverage} bounds the displacement a weighting induces by
the product of a leverage term, measuring how concentrated the weights are, and
a populational dispersion term. Reporting it alongside every comparison
separates the uninteresting case of an inert mechanism from the finding we
report: a mechanism that moves the search substantially and whose control
reproduces it exactly.

\subsection*{Contributions}

\begin{enumerate}
\item \textbf{Two controls for score attribution}, defined for a general
  mixing-operator form that subsumes cellular evolutionary algorithms,
  neighbourhood-topology swarms, localized consensus dynamics and
  centrality-weighted methods. They are drop-in and budget-neutral.

\item \textbf{Invariance results that make them sound.}
  Proposition~\ref{prop:mean} isolates assignment from heterogeneity---the
  confound that renders ablation inconclusive---and
  Proposition~\ref{prop:leverage} gives a necessary condition for a comparison
  to be informative at all, identifying the topologies on which such mechanisms
  are inert by construction.

\item \textbf{A large-scale negative result.} Across ten functions of the
  classic suite at two budgets and $28$ of CEC2017 at $D\in\{10,30,50\}$, more
  than $2\times10^{4}$ runs, betweenness-weighted learning is indistinguishable
  from its permutation control on every one of the $104$ function comparisons,
  while displacing each update by a third of the population's spread.

\item \textbf{Calibration across five outcomes.} The same controls, unchanged
  in suite, sample size and statistic, are applied to four further mechanisms
  across two published optimizers and return: inert by construction,
  uninformative, sign-varying by problem class, uniformly harmful, and
  uniformly beneficial. A null from a test with that range is a property of the
  mechanism.

\item \textbf{A reporting result.} On the fully informed particle swarm, a
  published quality weighting produces per-function effects reaching
  $|\delta|=1$ in both directions, significant on $15$ of $28$ functions, which
  cancel in the suite mean; a suite average alone would call this and the
  centrality null the same thing. The mirror failure appears in an adaptive
  scheme of our own, where per-function correction hid an effect that is
  unambiguous in aggregate. Both are demonstrated on data rather than argued in
  principle, and we subject our own proposal to the same control we ask of
  others.
\end{enumerate}

We do not claim that the optimizers studied perform poorly, or that score
weighting is useless in principle. The claim is narrower: for the mechanisms
tested here, the evidence offered for the attribution does not survive a
control that leaves the mechanism running. Whether other statistics, other
graphs or landscape-derived topologies fare differently is open, and the same
instrument settles it.
